\section{Evaluation}
We measure four cells at $\lambda=80$ (below the $128$-bit deployment level, the largest that fits the $24\,\mathrm{GB}$ budget; Table~1). Cells~1--3 share one scheme, so with/without/control is a matched comparison; Cell~4 (Aurora~\cite{cryptoeprint:2018/828}) is a different scheme over a different field, an indicative cross-scheme reference whose crossover is bounded in form, not located in wall-clock.

{\par\medskip\centering\footnotesize
\textbf{Table 1.}~Four cells (PLUM, $\lambda{=}80$, $24$\,GB, $n{=}5$ for cells~2--3).\\[3pt]
\begin{tabular}{cll}
\hline
Cell & Setting & Prove \\
\hline
1 & no precompile      & $>5$\,h \\
2 & Griffin precompile & $14.24$\,min \\
3 & SHA-3 control      & $13.24$\,min \\
4 & Aurora (Loquat)    & ${\approx}22$\,min \\
\hline
\end{tabular}\par\medskip}

With the precompile, verification fits the $24\,\mathrm{GB}$ budget for the first time in the study (peak ${\approx}14\,\mathrm{GB}$; Cell~1 does not terminate). The control is the sharpest result: even accelerated, PLUM verification proves $1.08\times$ slower and runs $1.13\times$ more cycles than the software SHA-3 (Table~1), so a standardized hash, unaccelerated, still beats our algebraic hash. The trace-area criterion postdicts the sign of this inversion, and a matched-field keystone ($164\times$ at $\ell=1$ versus $\ell=7$ limbs, a field-aligned operand versus the emulated $199$-bit one) fixes its necessity direction, so the inversion is evidence for the criterion, not a paradox.
